\documentclass{article}
\title{確率過程入門 — ランダムウォークから価格モデルへ}
\author{TeX 図書館}

\begin{document}

\section{この教科書のために — 「経路」を考える確率}

これまでの確率は「1 回の試行」や「1 つの数値の分布」でした。確率過程は\textbf{時間とともに変化する偶然}を扱います。サイコロの 1 目ではなく、サイコロを振り続けた\textbf{軌跡}。株価の明日の値ではなく、\textbf{価格の道すがら}。

\begin{quote}
\textbf{【なぜ経路が大事か】}\ 「期待値はプラス」と「資金が破綻しない」は別の問題でした(『投資の数学』第 4 節)。この差は\textbf{途中の経路}で決まります。経路の数学なしに破産確率は語れない、というのがこの教科書の出発点です。
\end{quote}

\section{ランダムウォーク — 最も単純な確率過程}

時刻 \(n\) の位置 \(X_n\) が「毎回 ±1 を等確率で足す」で決まる過程を考えます:
\[ X_{n+1} = X_n + \varepsilon_n, \qquad \varepsilon_n = \pm 1 \ (\text{各 } 1/2) \]

性質を数えるのは初等的です:
\begin{itemize}
\item \(E[X_n] = 0\)(期待値は動かない)
\item \(V[X_n] = n\)(ばらつきは \(\sqrt{n}\) で広がる)
\item 原点に戻る回数は無限に多く\textbf{なる}(1 次元のランダムウォークは再帰的)
\end{itemize}

\begin{quote}
\textbf{【例】}\ \(V[X_n] = n\) は「距離が \( \sqrt{n} \) で広がる」ことを意味します。10,000 ステップで標準偏差は 100 ステップぶん。緩やかですが確実に範囲が広がり、\textbf{戻る時間も伸びる}のが現実の相場や待ち行列と似ています。
\end{quote}

\begin{quote}
\textbf{【演習】}\ 勝率 \(p \neq 0.5\) のランダムウォーク(賭けの損益)では期待値のドリフトが \(E[X_n] = n(2p - 1)\) になることを示せ。
\end{quote}

\textbf{解答の指針:}\ 各ステップの期待値は \(1 \cdot p + (-1)(1-p) = 2p - 1\)。独立なので \(n\) 倍。ドリフトが正でも経路は大きく揺れる——期待値と経路の区別の第一歩。

\section{マルコフ連鎖 — 無記憶の偶然}

「次の状態が\textbf{現在の状態だけで}決まる」確率過程を\textbf{マルコフ連鎖}といいます。天気(今日の天気から明日の予報)や、ボードゲームの盤面が典型です。

遷移確率を表にまとめたものが\textbf{遷移行列} \(P\) です。たとえば「晴れ/雨」の 2 状態:
\[ P = \begin{pmatrix} 0.8 & 0.2 \\ 0.6 & 0.4 \end{pmatrix} \quad (\text{行は「今日の天気」、列は「明日の天気」}) \]
行の合計は 1(どこかへ遷移する確率の分配)です。

\(n\) 日後の天気分布は \(P\) を \(n\) 回掛ければ得られます:\(\boldsymbol{\pi}_n = \boldsymbol{\pi}_0 P^n\)。\textbf{遷移確率の繰り返し掛け算}こそが確率過程の計算の基本動作です。

\begin{quote}
\textbf{【演習】}\ 上の \(P\) で、今日晴れのとき明後日雨の確率を求めよ。
\end{quote}

\textbf{解答の指針:}\ \(P^2\) の (1,2) 成分 \(= 0.8 \times 0.2 + 0.2 \times 0.4 = 0.24\)。

\section{定常分布 — 動き続けても収束する先}

マルコフ連鎖を長く回すと、天気の予報分布がある分布 \(\pi\) に落ち着くことがあります:
\[ \pi P = \pi \]
この \(\pi\) を\textbf{定常分布}といいます。「動き続けているのに分布は安定」——経路は動いても、\textbf{問い合わせた時点での状態の割合}は一定になるということです。

\begin{quote}
\textbf{【例】}\ 上の \(P\) の定常分布を解く。\(\pi = (\pi_1, \pi_2)\) について \(\pi P = \pi\) と \(\pi_1 + \pi_2 = 1\) から
\[ 0.8\pi_1 + 0.6\pi_2 = \pi_1 \ \Rightarrow\ 0.6\pi_2 = 0.2\pi_1 \ \Rightarrow\ \pi_1 : \pi_2 = 3 : 1 \]
つまり長期的には「75\% 晴れ、25\% 雨」。\textbf{今日の天気の記憶は長期には消える}。
\end{quote}

線形代数の観点では、\(\pi P = \pi\) は「固有値 1 の固有ベクトルを求める」問題そのものです(『線形代数 応用編』第 9 節と直結)。

\begin{quote}
\textbf{【演習】}\ 毎回必ず反対に移る行列 \(P = \begin{pmatrix} 0 & 1 \\ 1 & 0 \end{pmatrix}\) の定常分布を求め、なぜ「落ち着かない」のか考察せよ。
\end{quote}

\textbf{解答の指針:}\ 解は \((0.5, 0.5)\) だが、実際の連鎖は晴れ↔雨を交互に繰り返すだけ(周期性)。\(\pi P = \pi\) の解と「実際の長期頻度」が一致するには\textbf{非周期的・連結}などの条件が要る。

\section{吸収確率 — ギャンブラーの破産を数える}

「0 に行ったら終わり(吸収)」という状態をもつランダムウォークで、破産する確率を厳密に計算できます。資金 \(i\)、勝率 \(p\)、負け率 \(q = 1-p\)、目標金額 \(N\) のとき、資金 \(i\) から\textbf{0 に先に到達する確率}は
\[ u_i = \begin{cases} 1 - \dfrac{1 - (q/p)^i}{1 - (q/p)^N} & (p \neq q) \\[2mm] 1 - \dfrac{i}{N} & (p = q) \end{cases} \]

\begin{quote}
\textbf{【例】}\ \(p = 0.5\)(フェア)、資金 \(i = 10\)、目標 \(N = 100\) なら破産確率は \(1 - 10/100 = 90\%\)! フェアゲームでも、\textbf{十分な利益額を目指すほど破産確率が上がる}——有限の資金で長く戦うほど運の順番に支配されるからです。\(p < 0.5\) なら(\(q/p)^N \to \infty\) で破産確率はほぼ 100\%。
\end{quote}

\begin{quote}
\textbf{【演習】}\ \(p = 0.51\), \(i = 10\), \(N = 20\) で破産確率を概算せよ(計算機可)。
\end{quote}

\textbf{解答の指針:}\ \(q/p = 0.49/0.51 \approx 0.9608\)。\(u_{10} = 1 - \frac{1 - 0.9608^{10}}{1 - 0.9608^{20}} \approx 1 - \frac{0.329}{0.549} \approx 40\%\)。わずかな優位でも、目標を近くすれば破産確率は半減する——\textbf{利益目標の設定がリスク}という教訓。

\section{ブラウン運動 — 極限としての連続ランダムウォーク}

ランダムウォークのステップ幅を \(\Delta x\)、時間間隔を \(\Delta t\) として \(\Delta x = \sqrt{\Delta t}\) の割合で細かくしていくと、連続的でギザギザな経路\textbf{ブラウン運動}(ウィーナー過程)が得られます。

性質は重要です:
\begin{itemize}
\item どの時点でも正規分布に従う(分散は経過時間に比例)
\item \textbf{どこでも微分できない}(どれだけ拡大してもギザギザ)
\item 変動幅は \(\sqrt{t}\) スケール
\end{itemize}

花粉の水中運動(ブラウン)から金融の価格モデルまで、連続的な偶然の標準モデルです。

\begin{quote}
\textbf{【演習】}\ ブラウン運動の「変動幅が \(\sqrt{t}\)」と、年間ボラティリティが 20\% の資産の「月間の変動目安は約 \(20\%/\sqrt{12} \approx 5.8\%\)」を結び付けて説明せよ。
\end{quote}

\textbf{解答の指針:}\ 期間 \(t\) の変動は \(\sigma\sqrt{t}\)。12 か月で 20\% なら 1 か月は \(20\% \times \sqrt{1/12}\)。ボラティリティのスケーリングは \(\sqrt{\text{時間}}\) が標準慣行。

\section{二項モデル — 離散で考える価格の未来}

株価の最も簡単なモデルが\textbf{二項モデル}です。各期間で価格が上がり \(u\) 倍、下がり \(d\) 倍の 2 択:

\[ S_0 \to \begin{cases} S_0 u \\ S_0 d \end{cases} \to \begin{cases} S_0 u^2,\ S_0 ud,\ S_0 d^2 \end{cases} \to \cdots \]

\(n\) 期後の価格は「上がった回数 \(k\)」だけで決まり、経路は \(\binom{n}{k}\) 通り。オプション価格は\textbf{各終点でのペイオフを確率で重みづけた期待値}として数え上げられます。

\begin{quote}
\textbf{【例】}\ \(S_0 = 100\), \(u = 1.2\), \(d = 0.9\), 2 期後の満期で行使価格 110 のコール。終点は \(144, 108, 81\)。ペイオフは \(34, 0, 0\)。中立確率 \(p^*\)(リスク中立確率)で期待値を取り、割り引いたものが理論価格。経路が見えるので、\textbf{「なぜこの価格か」が数え上げで説明できる}のが二項モデルの教育的価値です。
\end{quote}

\begin{quote}
\textbf{【演習】}\ 上の例でリスク中立確率 \(p^* = \frac{(1+r) - d}{u - d}\)(\(r\) は 1 期の無リスク利得)を、\(r = 0.01\) として計算し、コール価格の骨格を書け。
\end{quote}

\textbf{解答の指針:}\ \(p^* = \frac{1.01 - 0.9}{1.2 - 0.9} \approx 0.367\)。価格 \(= \frac{p^{*2} \times 34}{(1+r)^2} \approx \frac{0.1347 \times 34}{1.0201} \approx 4.5\)。

\section{モンテカルロ再訪 — 経路を乱数で歩く}

解析が難しい過程(可変ボラティリティ、複雑なルール)では、\textbf{経路そのものを乱数で大量生成}して分布を直接見るのがモンテカルロ法でした(『投資の数学』第 8 節)。確率過程の観点で言い直すと:

\begin{enumerate}
\item 経路を生む「ルール」(過程のモデル)を明文化する
\item 乱数で経路を \(N\) 本生成する
\item 経路から必要な統計量(最大下落幅・破産有無・満期ペイオフ)を計算する
\item その分布・タイル値で意思決定する
\end{enumerate}

\begin{lstlisting}[language=Python]
import random, statistics

def paths(f=0.1, p=0.55, b=1.0, n=200, trials=5000):
    finals, maxdds = [], []
    for _ in range(trials):
        eq, peak, dd = 1.0, 1.0, 0.0
        for _ in range(n):
            eq *= 1 + b * f if random.random() < p else 1 - f
            peak = max(peak, eq)
            dd = max(dd, 1 - eq / peak)
        finals.append(eq); maxdds.append(dd)
    return finals, maxdds

f_, d_ = paths()
print("最終資金の中央値:", round(statistics.median(f_), 3))
print("最大DDの90%タイル:", round(statistics.quantiles(d_, n=10)[-1], 3))
\end{lstlisting}

\begin{quote}
\textbf{【演習】}\ 上のシミュレーションで「最大ドローダウンの 90\% タイル」が 30\% を超えるなら、賭け割合をどう調整すべきか、ケリー基準の観点から述べよ。
\end{quote}

\textbf{解答の指針:}\ \(f\) をケリー比の半分以下に下げて再シミュレーションし、DD が許容範囲に収まる水準を探す。経路依存のリスクは「解析解でなく経路の分布」でしか設計できない。

\section{おわりに — 経路の目を手に入れる}

この教科書で手に入れた道具を振り返ります。ランダムウォーク(\( \sqrt{n} \) の広がり)、マルコフ連鎖(遷移行列と定常分布)、吸収確率(破産の厳密式)、ブラウン運動(連続の極限)、二項モデル(価格の数え上げ)、そしてモンテカルロ(経路の生成)。共通する見方は「\textbf{一回の損益でなく、経路の分布}」です。

『投資の数学』の資金管理、『統計学入門』の中心極限定理、『線形代数 応用編』のマルコフ連鎖の線形代数とつなげて読むと、確率過程は「偶然の時間構造」を統一する言葉であることが分かります。

\begin{quote}
\textbf{【総合演習】}
\begin{enumerate}
\item 勝率 0.6・ペイアウト比 1 のランダムウォークで、期待ドリフトと 100 ステップ後の標準偏差を求めよ。
\item 遷移行列 \(P = \begin{pmatrix} 0.5 & 0.5 \\ 0.25 & 0.75 \end{pmatrix}\) の定常分布を求めよ。
\item フェアゲーム(\(p = 0.5\))で資金 10、目標 1000 の破産確率を吸収確率の式で述べよ。
\end{enumerate}
\end{quote}

\textbf{解答の指針:}\ 1. ドリフト \(= 0.2/\)ステップ、標準偏差 \(= \sqrt{100} = 10\) ステップぶん。2. \(\pi_1 : \pi_2 = 1 : 2\)、つまり \((1/3, 2/3)\)。3. \(u_{10} = 1 - 10/1000 = 99\%\)。\textbf{フェアでも大目標はほぼ確実に破産}——資金管理の数学的根拠。

\end{document}
